\documentclass[11pt,a4paper]{article}

\usepackage[T1]{fontenc}
\usepackage{lmodern}
\usepackage{microtype}
\usepackage{amsmath,amssymb}
\usepackage{graphicx}
\usepackage{booktabs}
\usepackage[margin=2.5cm]{geometry}
\usepackage[hidelinks]{hyperref}

\title{{{title}}}
\author{{{author}}}
\date{\today}

\begin{document}
\maketitle

\begin{abstract}
  A short summary of the article.
\end{abstract}

\section{Introduction}
\label{sec:intro}
Write your text here. Cite sources like this~\cite{knuth1984}.
Section~\ref{sec:method} describes the method.

\section{Method}
\label{sec:method}
An equation with a number:
\begin{equation}
  E = mc^2 \label{eq:energy}
\end{equation}
Equation~\eqref{eq:energy} is famous.

\begin{figure}[htbp]
  \centering
  \fbox{\rule{0pt}{3cm}\rule{0.6\linewidth}{0pt}}
  \caption{Replace this box with \texttt{\textbackslash includegraphics}.}
  \label{fig:placeholder}
\end{figure}

\begin{table}[htbp]
  \centering
  \begin{tabular}{lrr}
    \toprule
    Item & Value & Unit \\
    \midrule
    Length & 1.5 & m \\
    Mass   & 3.2 & kg \\
    \bottomrule
  \end{tabular}
  \caption{A table with \texttt{booktabs}.}
  \label{tab:values}
\end{table}

\section{Conclusion}
Summarise your results.

\bibliographystyle{plain}
\bibliography{refs}

\end{document}
